\chapter{Gestion de projet}
\label{chap:gestion}

\iffalse
\section{Méthodologie}

Le projet a été conduit de façon incrémentale, organisée en sprints assortis de critères de
priorité (P0 indispensable, P1 important). Chaque sprint livrait un incrément vérifiable et
faisait l'objet d'une validation avant passage au suivant : modélisation Merise et UML,
puis implémentation, puis tests automatisés et audits de code. Cette discipline a permis de
détecter et corriger tôt les erreurs de modélisation (cardinalités, périmètre des entités).

\section{Organisation des sprints}

\begin{table}[H]
\centering
\caption{Découpage en sprints}
\label{tab:sprints}
\renewcommand{\arraystretch}{1.25}\footnotesize
\begin{tabularx}{\linewidth}{L{2.4cm} X}
\toprule
\textbf{Sprint} & \textbf{Contenu} \\
\midrule
Sprint 1 & Modélisation (MCD/MLD/UML), entités JPA, repositories, services de base, migrations Liquibase \\
Sprint 2 & Sécurité JWT + RBAC, DTO, pagination, gestion globale des erreurs, devises \\
Sprint 3 & Multi-tenant (TenantContext / TenantFilter), entités Message/Notification/Audit, analytique \\
Sprint 4 & Frontend Next.js (vitrine + tableau de bord), documentation, déploiement (Render, Vercel) \\
\bottomrule
\end{tabularx}
\end{table}

\section{Répartition des tâches et participation}

Le tableau~\ref{tab:participation} présente la répartition de la contribution entre les
membres du groupe. \emph{(Les taux sont à compléter ; chacun est évalué sur 100.)}

\begin{table}[H]
\centering
\caption{Répartition de la participation}
\label{tab:participation}
\renewcommand{\arraystretch}{1.3}
\begin{tabularx}{\linewidth}{L{6cm} X C{2.6cm}}
\toprule
\textbf{Membre} & \textbf{Contributions principales} & \textbf{Participation} \\
\midrule
ETCHOME Chris           & \dotfill & \underline{\hspace{1cm}} / 100 \\
{[Membre 2 — Nom Prénom]} & \dotfill & \underline{\hspace{1cm}} / 100 \\
KENNE DIHA Lea & \dotfill & \underline{\hspace{1cm}} / 100 \\
{[Membre 4 — Nom Prénom]} & \dotfill & \underline{\hspace{1cm}} / 100 \\
MBA BELA Oceanne & \dotfill & \underline{\hspace{1cm}} / 100 \\
{[Membre 6 — Nom Prénom]} & \dotfill & \underline{\hspace{1cm}} / 100 \\
\bottomrule
\end{tabularx}
\end{table}

\section{Outils et gestion de version}

Le code est versionné sous Git (dépôt GitHub), avec un flux de branches distinguant le code
stable (\texttt{main}), l'intégration (\texttt{dev}) et les branches de fonctionnalités
(\texttt{feat/*}) ou de correction (\texttt{fix/*}). La documentation technique est rédigée
en Markdown et en \LaTeX{} (le présent rapport).
\fi
